\begin{solution}
空气劈尖在反射光中有一次半波损失。距离劈棱 \(x\) 处膜厚近似为 \(d=x\theta\)，反射明纹满足总光程差为整数倍波长。

从劈棱起数，棱边处为暗纹，第四条明纹对应
 \(2x_4\theta+\frac{\lambda}{2}=4\lambda.\) 
于是
 \(x_4=\frac{7\lambda}{4\theta},\qquad x_4'=\frac{7\lambda}{4\theta'}.\) 
位移按新位置减旧位置记：
 \(\Delta x=x_4'-x_4 =\frac{7\lambda}{4}\left(\frac{1}{\theta'}-\frac{1}{\theta}\right) =\frac{7\lambda(\theta-\theta')}{4\theta\theta'}.\) 
第四条明纹离棱边不是 \(4\) 个条纹间距，而是 \(3.5\) 个条纹间距，因为反射光在棱边处为暗纹。若 \(\theta'\) 变小，条纹间距变大，第四明纹应向远离棱边方向移动，公式也给出正位移。
\end{solution}


\begin{exercise}[用劈尖条纹间距测金属细丝直径时为什么先写几何斜率]
把金属细丝夹在两块平玻璃之间形成空气劈尖。用波长 \(\lambda=\SI{589.3}{nm}\) 的单色光垂直照射，已知金属丝到劈尖顶点的距离 \(L=\SI{28.880}{mm}\)，测得 \(30\) 条明纹跨越的总长度为 \(\SI{4.295}{mm}\)。求金属细丝直径 \(D\)。
\end{exercise}

\begin{solution}
这类题看上去在测长度，其实真正要写的是“厚度沿长度方向怎样线性增长”。也就是说，先把劈尖的几何斜率写成 \(\alpha\approx D/L\)，再用条纹间距把 \(\alpha\) 反推出去。

这里空气劈尖取 \(n=1\)，且 \(30\) 条明纹跨越的是 \(29\) 个条纹间距，所以
\[
l=\frac{4.295}{29}\,\si{mm}\approx \SI{0.1481}{mm},\qquad
D=L\alpha=L\frac{\lambda}{2l}
=\frac{28.880\times 589.3\times10^{-6}}{2\times0.1481}\,\si{mm}
\approx \SI{0.0574}{mm}.
\]
细丝直径约为 \(\SI{5.74e-2}{mm}=\SI{57.4}{\micro m}\)，落在几十微米量级是合理的。这里不能把 \(30\) 条条纹误当成 \(30\) 个间距，也不能漏掉空气层对应 \(n=1\)。
\end{solution}


\begin{exercise}[空气劈尖条纹弯向厚端时怎样判断凸凹缺陷]
光学平板玻璃 \(A\) 与待测工件 \(B\) 间形成空气劈尖，左侧为空气膜薄端、右侧为空气膜厚端。用波长 \(\lambda=\SI{500}{nm}\) 的单色光垂直照射，反射光干涉条纹中弯曲部分向厚端偏移，且顶点恰好与其右边相邻直条纹的连线相切。判断工件上表面缺陷是凸起还是凹槽，并求最大高度或深度。
\end{exercise}

\begin{solution}
题目说弯曲顶点恰好与右边相邻直条纹相切，表示缺陷导致的最大膜厚变化正好相当于相邻一级条纹，即
 \(h_{\max}=\frac{\lambda}{2} =\frac{\SI{500}{nm}}{2} =\SI{250}{nm}.\) 
所以缺陷为凸起，最大高度为 \(\SI{250}{nm}\)。
\end{solution}
